% section_4_additional.tex — v6 from-scratch rewrite (Phase C 4-col rebuild)
% Spec: docs/Draft/THESIS_SKELETON.md v6 §4 outline
% §4.1 = drivers (Phase 5); §4.2 = outsider reaction (Phase 6); §4.3 = endogeneity (Phase 7)
% Numbers: programmatic extract from suite_specs (post-Phase-C 4-col CEO-only DVs)
%   H11 PRisk:        n 44,497–65,760  (UncResCEO + UncPreCEO × 2 FE = 4 cells)
%   H11-Lag PRisk_l2: n 41,837–60,503
%   H23 TSIMM:        n 12,061–18,492
%   H24 US EPU:       n 39,063–60,503
%   H24b GEPU:        n 39,063–60,503

\section{Additional Analyses}
\label{sec:additional}

\subsection{Drivers of CEO Speech Uncertainty}
\label{sec:additional:drivers}

The cash-financing and outside-world results in Sections~\ref{sec:main} and \ref{sec:additional:reaction} treat the speech-uncertainty constructs as right-hand-side regressors. Their interpretation as legitimate uncertainty proxies depends on whether they respond, on the left-hand side, to independently measured exogenous uncertainty drivers. We test four such drivers as construct-validation regressions of the speech-uncertainty constructs on (i) the firm-level political-risk index of \citeA{hassan2020}, (ii) the U.S.\ news-based economic-policy-uncertainty index of \citeA{baker2016}, (iii) the global extension of \citeA{davis2016}, and (iv) product-market competition pressure as measured by the text-based total similarity index of \citeA{hoberg2016}. Each driver regression is reported across four cells: the two CEO speech-uncertainty components from the \citeauthor{dzielinski2021} decomposition (UncResCEO Q\&A residual and UncPreCEO presentation segment) crossed with two entity fixed-effect configurations (industry vs.\ firm). The macro-driver suites (BBD-EPU and GEPU) use two-way clustering on firm and calendar quarter to handle within-quarter cross-firm correlation in the macro regressor, per Section~\ref{sec:framework:design}. The four drivers test construct validity rather than formal hypotheses.

\textbf{Political risk (Hassan et al., 2020).} The contemporaneous firm-level political-risk index loads positively on both CEO speech-uncertainty components in all four cells of Table~\ref{tab:h11_prisk_uncertainty} at $p<0.01$ one-tailed, with point estimates between $+0.0001$ and $+0.0003$. The two-quarter-lag specification (Table~\ref{tab:h11_lag2}) is positive in all four cells, with the two UncPreCEO cells significant at $p<0.01$ ($\beta\approx +1.5\times 10^{-4}$ industry FE, $\beta\approx +3.7\times 10^{-5}$ firm FE); the two UncResCEO cells are positive but non-significant at $p<0.10$ (point estimates near zero, $p_{\text{one}}\approx 0.18$--$0.20$). The lag specification establishes that the firm-level political-risk signal is a leading driver of the presentation-segment uncertainty share, not merely a contemporaneous correlation; the Q\&A-residual share carries no two-quarter-ahead political-risk signal.

\textbf{Macro economic-policy uncertainty (Baker, Bloom, and Davis, 2016; Davis, 2016).} The U.S.\ economic-policy-uncertainty index (Table~\ref{tab:h24_us_epu}) loads positively on both CEO speech-uncertainty components in all four cells, with all four significant at $p<0.10$ (three of four at $p<0.05$ and two of four at $p<0.01$), point estimates between $+0.0123$ and $+0.0239$. The global extension (Table~\ref{tab:h24b_global_epu}) is similarly positive in all four cells, with all four significant at $p<0.05$ (two of four at $p<0.01$), point estimates between $+0.0181$ and $+0.0309$. The two macro indices both deliver the predicted construct-validation pattern: the CEO speech-uncertainty components covary with independently measured macro uncertainty.

\textbf{Product-market competition (Hoberg and Phillips, 2016).} The text-based total-similarity index loads positively in all four cells of Table~\ref{tab:h23_competition_uncertainty} but is concentrated in the presentation-segment component: both UncPreCEO cells are significant at $p<0.01$ ($\beta = +0.0304$ industry FE, $\beta = +0.0302$ firm FE), while the two UncResCEO cells are positive but non-significant at $p<0.10$ ($\beta = +0.0008$ industry FE with $p_{\text{one}}=0.29$, $\beta = +0.0023$ firm FE with $p_{\text{one}}=0.38$). The Pres-channel concentration is consistent with \citeauthor{dzielinski2021}'s framing of the prepared-remarks segment as the firm-culture-reflecting share: competition pressure manifests in the segment of the call constructed in advance with the firm's broader competitive context in mind, while the Q\&A residual carries no detectable competition signal.

\textbf{Synthesis.} All four uncertainty drivers load positively in every cell, with $p<0.10$ significance covering 4 of 4 (political risk), 4 of 4 (US-EPU), 4 of 4 (GEPU), and 2 of 4 (competition) cells. The two non-significant competition cells are both UncResCEO cells ($p_{\text{one}}>0.28$); none of the sixteen driver--component cells reverses sign. The political-risk lag specification confirms timing for the presentation-segment component. Sign and significance evidence across the four drivers establishes the two CEO speech-uncertainty components as legitimate uncertainty proxies before the cash-financing and outside-world outcome tests in Sections~\ref{sec:main} and \ref{sec:additional:reaction}.

\subsection{Outsider Reaction: Spread and SEC Comment Letter}
\label{sec:additional:reaction}

Section~\ref{sec:main} reports the firm-side cash-financing response to CEO speech uncertainty under the three-component decomposition. We now turn to two outsider-reaction outcomes that \citeauthor{dzielinski2021} do not test, both at the firm-quarter resolution: the post-call quoted bid-ask spread (a market-microstructure outcome attributable to outside traders' information-processing costs) and the receipt of an SEC comment letter referencing the earnings call (a regulatory outcome attributable to SEC reviewers' attention to call content). The two outsider-reaction tests adopt the same three-component decomposition as the firm-side cash test, jointly entered as right-hand-side regressors with Bates-style controls. The two tests extend the DWZ measurement architecture to two new outcome classes: market microstructure beyond the abnormal-volume metric \citeauthor{dzielinski2021} use, and regulatory-reaction outcomes their paper does not address.

\subsubsection{Market Channel: Post-Call Bid-Ask Spread}
\label{sec:additional:reaction:spread}

The post-call quoted bid-ask spread aggregates outside traders' information-asymmetry response to call content over the post-disclosure information-incorporation window. \citeA{bushee2018} construct an information-asymmetry measure as the average value of the \citeA{amihud2002} illiquidity construct ``over the period starting the day of the call and ending 25 trading days subsequent to the call.'' Because \citeauthor{bushee2018}'s alternative spread-side measure uses an intraday TAQ-based structural estimator that is not feasible in our daily-CRSP setting, we adapt their post-call 25-trading-day window definition (day~0 inclusive, per \citeauthor{bushee2018} verbatim) to a daily closing-quote relative bid-ask spread averaged over the window. Both the \citeauthor{bushee2018} window and the daily-CRSP-feasible spread formula are documented in Appendix~\ref{app:vardefs}.

Table~\ref{tab:h14c_ceo2_decomp} reports the four-step fixed-effects ladder of post-call-spread regressions on the three speech-uncertainty components. The presentation share \texttt{UncPreCEO} loads positively in 4 of 12 specifications at $p<0.10$, with 3 of 12 at $p<0.05$ and 2 of 12 at $p<0.01$. Per-cell inspection of the suite specification confirms the four significant cells are uniformly positive: cols 1 ($\beta=+0.1664$, $p_{\text{one}}=0.0393$), 2 ($\beta=+0.3496$, $p_{\text{one}}=0.0026$), 4 ($\beta=+0.2945$, $p_{\text{one}}=0.0083$), and 6 ($\beta=+0.1644$, $p_{\text{one}}=0.0825$), all on the contemporaneous post-call-spread dependent variable. The eight cells not statistically distinguishable from zero include four lead-DV cells with negative point estimates, all of them non-significant at $p<0.10$. The presentation-segment uncertainty share predicts the call-quarter post-call spread but does not predict the next-quarter spread; the effect is contemporaneous-only, consistent with information-processing costs absorbed within the BGT post-disclosure window.

The Q\&A components are null. \texttt{ClarityCEO} is not statistically distinguishable from zero in 0 of 12 cells at $p<0.10$, with point estimates spanning $-0.1886$ to $+1.1056$ across cells (both signs across cells, all non-significant). \texttt{UncResCEO} is null in 0 of 12 cells at $p<0.10$, with point estimates between $-0.1021$ and $+0.0984$; the contemporaneous cells uniformly negative ($-0.10$ to $-0.06$) and the lead cells uniformly positive ($+0.05$ to $+0.10$), but neither margin distinguishable from zero.

\subsubsection{Regulatory Channel: SEC Comment Letter Receipt}
\label{sec:additional:reaction:sec}

Following \citeA{lerman2026}, we measure SEC scrutiny attributable to call content using the conference-call-comment-letter (CCCL) indicator: an indicator that the firm receives an SEC comment letter referencing the focal earnings call within the call's fiscal quarter. The indicator is rare in the panel (call-quarter receipt rate well below one percent), and we estimate it as a linear-probability model in PanelOLS to maintain comparability with the rest of the body display.

Table~\ref{tab:h18_ceo2_decomp} reports six specifications crossing two entity fixed-effect configurations (industry vs.\ firm) with two control specifications (base vs.\ extended) and two time-fixed-effect configurations (calendar year vs.\ calendar year-quarter). \texttt{UncPreCEO} loads positively on CCCL receipt in 4 of 6 cells at $p<0.10$ (cols 2, 3, 4, 6), with 1 of 6 at $p<0.05$ (col 3, $\beta=+0.0016$, $p_{\text{one}}=0.0144$). Per-cell inspection of the suite specification confirms all six cells are positive (point estimates between $+0.0007$ and $+0.0016$). Three of the four extended-controls specifications (cols 3, 4, 6) survive at $p<0.10$; col 5 (industry-FE, calendar year-quarter time effects) is the lone non-significant extended-controls cell at $p_{\text{one}}=0.1652$. The presentation-segment loading on regulatory-reaction outcomes survives the more demanding within-firm-and-within-quarter identifying variation in three of four extended-control configurations.

The Q\&A components are null on the regulatory-channel outcome, mirroring the spread-channel pattern. \texttt{ClarityCEO} is not statistically distinguishable from zero in 0 of 6 cells at $p<0.10$, with point estimates between $+0.0039$ and $+0.0062$ but all $p_{\text{one}}>0.50$. \texttt{UncResCEO} is null in 0 of 6 cells at $p<0.10$, with near-zero point estimates between $-0.0004$ and $-0.0003$. Neither Q\&A component carries marginal information for SEC comment-letter receipt once the presentation share is jointly included.

\subsubsection{Insider-Outsider Channel Asymmetry}
\label{sec:additional:reaction:asymmetry}

The Section~\ref{sec:main} firm-side cash response loads on the Q\&A components (\texttt{ClarityCEO} negative in 9 of 12 specifications, \texttt{UncResCEO} positive in 12 of 12); the Sections~\ref{sec:additional:reaction:spread} and \ref{sec:additional:reaction:sec} outsider-reaction outcomes load on the presentation-segment component (\texttt{UncPreCEO} positive in 4 of 12 spread cells and 4 of 6 SEC cells, with the Q\&A components null in both regressions). The asymmetry sits cleanly on \citeauthor{dzielinski2021}'s own theoretical structure for the two segments of the call. The presentation segment is the one \citeauthor{dzielinski2021} characterize as ``scripted and vetted beforehand\dots\ likely to reflect the firm's culture''; the Q\&A segment is the one they characterize as ``inevitably somewhat improvised and hence more likely to reveal each CEO's style.'' Two segments of the same call carry signals to actors with different information sets and different decision horizons. The thesis documents this insider-outsider channel asymmetry within the DWZ framework while extending the architecture's outcome set from market reactions, analyst behavior, performance, and governance to firm financing-policy and regulatory-reaction outcomes.

\subsection{Endogeneity}
\label{sec:additional:endo}

\subsubsection{Three Designs Targeting Three Orthogonal Threats}
\label{sec:additional:endo:framing}

The cash response of Section~\ref{sec:main} faces three identification threats. First, reverse causality: cash policy could shape CEO speech rather than the reverse, with cash-rich firms speaking with calmer or clearer cadence regardless of underlying uncertainty. Second, time-invariant heterogeneity: a CEO-firm match could embed both speech style and cash policy through the matching process, with neither causing the other within match. Third, time-varying omitted confounders: a within-CEO time-varying state variable (a strategy review, a litigation flare-up, a customer-concentration shift) could move both speech uncertainty and cash policy contemporaneously, producing the within-firm correlation Section~\ref{sec:main} documents.

We address these three threats with three designs that target them orthogonally rather than converge on a single point estimate. Subsection~\ref{sec:additional:endo:death} reports a CEO sudden-death difference-in-differences against the reverse-causality threat. Subsection~\ref{sec:additional:endo:dwzfd} reports a first-difference replication of \citeauthor{dzielinski2021}'s Section~6 turnover specification against the time-invariant heterogeneity threat. Subsection~\ref{sec:additional:endo:lewbel} reports a \citeA{lewbel2012} heteroskedasticity-based IV against the time-varying omitted-confounders threat. Subsection~\ref{sec:additional:endo:asymmetry} discloses a structural asymmetry between the ClarityCEO and UncResCEO identification problems and disposes of the asymmetric coverage of the three designs.

\subsubsection{CEO Sudden-Death Difference-in-Differences}
\label{sec:additional:endo:death}

Following \citeA{bennedsen2020}, who identify CEO causal effects using exogenous shocks to CEO availability (hospitalization and death), and \citeA{ghafoor2023}, who apply a sudden-death DiD to firm cash holdings, we construct an exogenous-shock test on our F1D panel. We classify CEO-death events using \citeauthor{ghafoor2023}'s operational definition of sudden death: heart attack, plane or automobile accident, or news language indicating the death was ``due to sudden illness'' or ``passed away unexpectedly,'' excluding deaths preceded by visible decline. Of the candidate events identified in the 2002--2018 panel, eight have the call-coverage and pre-event panel availability required for a $\pm$12-quarter DiD design with at least five pre-event Q\&A-segment calls per treated CEO.

Treated firms are matched 1:1 to control firms via nearest-neighbor propensity-score matching on \citeA{bates2009} covariates at $t-1$ (lnAssets, leverage, ROA, dividend dummy, scaled cash flow), drawn from the non-treated F1D pool within the same death-quarter $\pm$2 window. Post-matching standardized differences are within the Rosenbaum--Rubin threshold of 0.25 on all five covariates. Table~\ref{tab:ceo_death_did} reports four DiD specifications crossing two entity fixed-effect configurations (industry vs.\ firm) with two time fixed-effect configurations (none vs.\ year-quarter). The pooled-ATT estimate (\texttt{Treated}~$\times$~\texttt{Post}) is negative across all four specifications, ranging from $\beta=-0.0053$ (firm$+$year-quarter FE) to $\beta=-0.0178$ (industry FE only, replicating \citeauthor{ghafoor2023}'s baseline specification), with one-tailed p-values between $0.121$ and $0.338$ on a treated-firm count of eight. The negative sign matches \citeauthor{ghafoor2023}'s reduced-form magnitude and the precautionary-cash prior; the magnitudes are 12--40\% of \citeauthor{ghafoor2023}'s primary point estimate of $-0.043$. None of the four specifications is statistically significant at conventional one-tailed thresholds.

Three honest disclosures qualify the design at this sample size. \emph{Pre-trend pattern.} A pre-period placebo test on the industry-FE specification estimates Treated~$\times$~Pre$_{4q}$ at $\beta=+0.027$ ($p_{\text{two}}=0.183$) and Treated~$\times$~Pre$_{8q}$ at $\beta=+0.021$ ($p_{\text{two}}=0.146$); both are statistically insignificant but positive in sign while the post-event ATT is negative. The pattern is consistent with mean reversion rather than a clean parallel-trends pass; we describe it as such rather than claim ``parallel trends pass.'' \emph{Lagged-DV exclusion.} The Phase~E specifications omit the lagged dependent variable that is otherwise required in Section~\ref{sec:main}, following the standard DiD practice of \citeA{bertrand2004}: cash is highly persistent, and including a lagged dependent variable in the DiD specification mechanically absorbs the ATT through cash-stock persistence. The firm fixed effects in cols~2 and~4 already address time-invariant heterogeneity. \emph{Cluster count.} With sixteen unique firm clusters (eight treated plus eight matched controls), cluster-robust standard errors are below the conventional asymptotic-validity threshold of approximately thirty; we report the conventional cluster-robust standard errors with this disclosure. The originally proposed heterogeneity test on pre-event UncResCEO is not feasible at this sample size (a median split would yield four events per cell); we report the pooled ATT only and treat the speech-channel-specific decomposition as future work conditional on a larger event sample.

\subsubsection{First-Difference Replication of Dzielinski, Wagner, and Zeckhauser Section~6}
\label{sec:additional:endo:dwzfd}

\citeauthor{dzielinski2021} address time-invariant firm and manager heterogeneity in their Section~6 by estimating a first-difference specification around CEO-turnover events: the change in the firm-level outcome is regressed on the change in average ClarityCEO between the outgoing-CEO tenure and the incoming-CEO tenure, with sample filtered to turnovers where (i) the gap between the outgoing CEO's last call and the incoming CEO's first call is at most 120 days, (ii) both CEOs have at least five calls at the firm, and (iii) both appear in the underlying CEO fixed-effect estimation window. The first-difference annihilates all firm-level and manager-level fixed components; identification is from within-turnover-pair change.

We adapt the \citeauthor{dzielinski2021} Section~6 design from their Tobin's Q outcome to cash holdings. The sample-construction filters yield $659$ turnover pairs in the F1D panel. Table~\ref{tab:h_dwz_fd} reports three specifications: the verbatim \citeauthor{dzielinski2021} controls (col~1), the verbatim controls plus the additional Bates extras (col~2), and an industry-adjusted variant (col~3). The coefficient on $\Delta$ClarityCEO is $\beta=-0.0183$ ($p_{\text{one}}=0.046$) in col~1, $-0.0172$ ($p_{\text{one}}=0.054$) in col~2, and $-0.0182$ ($p_{\text{one}}=0.046$) in col~3, each negative and significant at $p<0.10$ one-tailed in the predicted direction (clarity rises after turnover $\Rightarrow$ cash falls, the within-firm-and-within-pair-of-CEOs analogue of the within-firm pattern in Section~\ref{sec:main:hc}). The Phase~E and DWZ-first-difference designs converge on identical $|\beta|\approx 0.018$ across an exogenous-low-power and a weak-identification-high-power contrast on the same outcome (cash holdings) and the same speech-component (ClarityCEO).

The first-difference does not identify the UncResCEO effect because the within-CEO mean of UncResCEO is zero by construction (the residual from a CEO-fixed-effect regression has within-CEO mean exactly zero by the OLS first-order conditions; verified empirically at machine zero on our panel). The change in the per-CEO mean of UncResCEO between two CEOs is therefore identically zero in expectation, and the first-difference specification identifies ClarityCEO only. Three deviations from the \citeauthor{dzielinski2021} Section~6 specification are documented in Appendix~\ref{app:vardefs}: (i) industry fixed effects use the Fama-French~12 partition rather than \citeauthor{dzielinski2021}'s Fama-French~48 (panel-infrastructure constraint); (ii) the intangibles-ratio control is omitted because it is not constructed in our panel; (iii) the sample restricts to F1D Main industries (Fama-French~12 excluding Finance and Utilities) per the rest of the body. Cash is also a CEO-discretion choice variable, which makes endogeneity more concerning here than in \citeauthor{dzielinski2021}'s Tobin's Q application. We treat this design as a robustness companion to the Phase~E sudden-death DiD rather than as primary identification.

\subsubsection{Heteroskedasticity-Based Instrumental Variable Estimation}
\label{sec:additional:endo:lewbel}

The \citeA{lewbel2012} estimator constructs instruments from heteroskedasticity in the residual of the endogenous regressor relative to a vector of base controls. Under the exclusion restriction that each base control has covariance with the squared residual of the endogenous regressor but is uncorrelated with the structural-equation error term, the heteroskedasticity-based instruments identify the structural coefficient on the endogenous regressor. The base controls in our application are the eight \citeauthor{bates2009} controls; we apply a heteroskedasticity-screening filter to retain only those controls whose squared residual has a statistically significant relation to the endogenous regressor's residual. Six of the eight \citeauthor{bates2009} controls are retained (lnAssets, leverage, Tobin's Q, ROA, capex, and the lagged cash ratio); dividend dummy and scaled cash flow are dropped at the screening stage.

Table~\ref{tab:h_lewbel_iv} reports three specifications. Col~1 estimates the OLS baseline as a replication check: the coefficient on UncResCEO is $\beta=+0.0021$ ($p_{\text{one}}=0.020$, $n=43{,}471$), which replicates the cell~1 estimate of the main panel in Section~\ref{sec:main:hc} (the small observation difference relative to the main panel's $n=43{,}333$ reflects control-NaN listwise drops in the IV-instrument residualization step). Col~2 reports the 2SLS-Lewbel estimate with the six-instrument heteroskedasticity-screened set: $\beta=+0.0100$ ($p_{\text{one}}=0.115$, $n=43{,}471$). Col~3 reports the 2SLS-Lewbel estimate with the extended nine-instrument set: $\beta=+0.0101$ ($p_{\text{one}}=0.116$, $n=43{,}454$).

The diagnostic battery is reported in the table footnote. The first-stage Cragg-Donald F-statistic is $20.4$ in col~2 and $41.3$ in col~3, both above the standard weak-IV threshold of $F=10$; the col~2 statistic is below the more demanding Stock-Yogo $10$\%-maximal-IV-size threshold of approximately $23$ for one endogenous regressor with six instruments, so we disclose the col~2 first-stage as borderline weak under the Stock-Yogo criterion. The Sargan over-identification test fails to reject the joint-validity null in col~2 ($p=0.92$) but rejects it at $p<0.05$ in col~3 ($p=0.009$); we accordingly treat col~2 as the primary specification and disclose that the extended-instrument col~3 fails the over-identification test, likely reflecting one of the additional extended-controls $W_j$ violating the heteroskedasticity-mean-independence exclusion restriction. The Wu-Hausman test of OLS-consistency under exogeneity does not reject the null in col~2 ($p=0.24$); we treat this as a failure to reject rather than evidence in favor of OLS-consistency. The 2SLS-Lewbel point estimate is approximately five times the OLS estimate, a magnitude consistent with classical measurement-error attenuation in OLS; the 2SLS-Lewbel and OLS point estimates are not statistically distinguishable from each other at conventional thresholds given the wider 2SLS standard error.

\subsubsection{ClarityCEO vs.\ UncResCEO: Identification Asymmetry}
\label{sec:additional:endo:asymmetry}

The three designs target different speech components by construction. The Phase~E and DWZ first-difference designs identify ClarityCEO via cross-CEO-pair variation around turnover events; neither identifies UncResCEO, because the residual of a CEO-fixed-effect regression has within-CEO mean exactly zero by OLS first-order conditions, so any first-difference that averages within-CEO has $\Delta$UncResCEO $\equiv 0$ in expectation. The Lewbel design identifies UncResCEO via call-level heteroskedasticity-based instruments; it does not separately identify ClarityCEO because ClarityCEO is a CEO-level trait (one observation per CEO in the panel) whose call-level variation is, by construction, the variation absorbed into UncResCEO. The two regressors therefore have structurally distinct identification problems and require structurally distinct identification tools. The asymmetry across designs reflects the underlying statistical structure of the two speech components, not a coverage gap in the endogeneity package.
